\section{Objectives and Research Questions}

The main objective of this thesis is to investigate the effectiveness of deep
learning-based approaches for automatic source code translation, with a
particular focus on the effect of fine-tuning and the reliability of different
evaluation methods. Although modern large language models already perform well
on many code-related tasks, it is still unclear to what extent fine-tuning can
improve translation quality in terms of both syntactic accuracy and functional
correctness. %szb: ez miért unclear? indokolni kéne, pl. kevés study van róluk, kevés benchmark, vagy ilyesmik - és persze beciteolni azokat, amik vannak, vagy mégjobb, beciteolni azt, hogy más unclear-nak nevezi ezt

To address this problem, the thesis compares base and fine-tuned versions of %szb: lehetne egyértelműség kedvéért kicsit explicitebb, hogy miben hasonlítod őket össze
multiple code-oriented language models under the same experimental conditions.
The models are fine-tuned on the Java-to-C++ portion of the XLCoST dataset %szb: cite lemaradt 
and evaluated on the CodeEditorBench benchmark %szb: cite lemaradt
using consistent preprocessing, prompting, and evaluation procedures. This setup
makes it possible to compare different model architectures as well as the impact %szb: itt nem kéne már picit explicitebbnek lenni? nem tudom...
of fine-tuning in a fair and reproducible way.

A further objective is to investigate the relationship between similarity-based
metrics and execution-based evaluation methods. Previous studies have shown that
metrics such as BLEU often correlate poorly with functional correctness,
motivating the use of execution-based evaluation \cite{evaluation_limitations,
huang2022exeds}. This thesis therefore examines both types of evaluation in
order to obtain a more complete picture of code translation quality. %szb: le kéne írni, hogyan történik az examination

Based on these objectives, the thesis is guided by the following research
questions:

%szb: amúgy most elgondolkoztam... az RQ1 meg 2 az nem ugyanaz?
% hisz egy translation akkor helyes, ha megőrzi a működést. A modell translation effectiveness-t pedig úgy mérjük, hogy ez hánszor történik
\begin{itemize}
    \item \textbf{RQ1: How effective are deep learning models in translating %szb: lehetne konkretizálni, hogy llm, hisz az van fókuszban
    source code between programming languages?}

    \item \textbf{RQ2: To what extent do these models preserve the functional %szb: nem tudom, hogy a preserve biztos jó szó-e... Esetleg transfer behaviour? vagy ilyesmi? am maradhat ha nincs jó ötlet
    correctness of the original code?} 

    \item \textbf{RQ3: How effective are execution-based evaluation methods
    compared to traditional similarity-based metrics?} %szb: nem jó sztem a szóhasználat: effective a modellekre használandó sztem, itt inkább az a kérdés, hogy:
    %a tradicionális metrikák mennyire jók, illetve hogyan viszonyulnak az execution-based metrikákhoz?
    %kb ez, nem? nem biztos

    \item \textbf{RQ4: What types of errors are most commonly produced by deep
    learning models during code translation?}
\end{itemize}
%ja, és hiányolom a finetuning aspektusát is a rq-kból
%kicsit szkeptikus vagyok ezekkel a kérdésekkel kapcsolatban, bocsána,t hogy csak most írom. javaslok alternatívákat:
% 1 - how effective are current small language models in code translation
% 2 - to what extent can fine tuning improve their translation abilities
% 3 - What types of errors are most commonly produced
% 4 - How well do traditional metrics (such as bleu) estimate the effectiveness of translation llms and how do they relate to execution-based evaluation